% Part: turing-machines
% Chapter: undecidability
% Section: universal-tm

\documentclass[../../../include/open-logic-section]{subfiles}

\begin{document}

\olfileid{tur}{und}{uni}
\olsection{सार्वत्रिक ट्यूरिंग यंत्रे}

\olref[enu]{sec} मध्ये प्रत्येक ट्यूरिंग यंत्राचे
पूर्णांकांच्या सांत क्रमिकेने वर्णन कसे करता येते,
याची चर्चा केली. ही क्रमिका त्या यंत्राच्या अवस्था,
वर्णमाला, आरंभीची अवस्था आणि सूचना सांकेतिक करते.
अशा सर्व वर्णनांचा संच !!{enumerable} म्हणजे गणनीय
आहे, हेही आपण दाखवले. अशा वर्णनांचा संच
!!{denumerable} म्हणजे गणनीय अनंत असल्यामुळे
~$\Nat$ पासून या वर्णनांवर जाणारे
!!a{surjective} म्हणजे आच्छादक फलन आहे.
असे !!a{surjective} फलन उदाहरणार्थ कांतोरची
नागमोडी पद्धत वापरून मिळवता येते. त्यामुळे
ट्यूरिंग यंत्रांच्या (वर्णनांचे) प्रगणन करता येते.
यापैकी एक प्रगणन निश्चित केल्यावर $1$-वे,
$2$-रे, \dots, $e$-वे ट्यूरिंग यंत्र असे
बोलता येते. या संख्यांना \emph{निर्देशांक} म्हणतात.

\begin{defn}
$M$ हे (आपण निश्चित केलेल्या प्रगणनातील)
$e$-वे ट्यूरिंग यंत्र असेल, तर $e$ हा~$M$ चा
\emph{निर्देशांक} आहे असे म्हणतो. $e$-व्या
ट्यूरिंग यंत्रासाठी $M_e$ असे लिहितो.
\end{defn}

एका यंत्राला एकाहून अधिक निर्देशांक असू शकतात.
उदाहरणार्थ, $M$~च्या सूचनांची यादी कोणत्या
क्रमाने करतो यावरून त्याची दोन वर्णने वेगळी
असू शकतात; अशा वेगळ्या वर्णनांचे निर्देशांकही
वेगवेगळे असतात.

महत्त्वाची बाब अशी की ट्यूरिंग यंत्रांच्या वर्णनांचे
प्रगणन अशा रीतीने करता येते की निर्देशांकावरून
यंत्र~$M$ चे वर्णन प्रभावी रीतीने संगणित करता
येईल आणि यंत्र~$M$ च्या वर्णनावरून त्याचा
एक निर्देशांकही प्रभावी रीतीने संगणित करता
येईल. मग चर्च--ट्यूरिंग प्रबंधानुसार, निर्देशांक~$e$
असलेल्या ट्यूरिंग यंत्राचे वर्णन मिळवून ते
आपल्या फितीवर फलित म्हणून लिहिणारे ट्यूरिंग
यंत्र असू शकते. हे वर्णन म्हणजे
~$\TMstroke$ चिन्हांच्या खंडांची क्रमिका असेल
(हे खंड $M_e$ चे वर्णन करणाऱ्या क्रमिकेतील
धन पूर्णांक दर्शवतील).

आता असा प्रश्न स्वाभाविकपणे पडतो: ट्यूरिंग
यंत्रांच्या निर्देशांकांवरील कोणती फलने स्वतः
ट्यूरिंग यंत्रांनी संगणित करता येतात? ट्यूरिंग
यंत्रांच्या निर्देशांकांचे कोणते गुणधर्म ट्यूरिंग
यंत्रांनी निर्णीत करता येतात? उदाहरण पाहा:
निर्देशांक~$e$ वरून निर्देशांक~$e$ असलेल्या
ट्यूरिंग यंत्रातील अवस्थांची संख्या देणारे
फलन ट्यूरिंग यंत्राने संगणित करता येते.
असे यंत्र पुढीलप्रमाणे काम करेल: सुरुवातीला
त्याच्या फितीवर $e$~$\TMstroke$ चिन्हांचा
एकच खंड असेल. प्रथम ते $e$ वरून त्याचे
वर्णन उलगडेल. आता ते वर्णन फितीवरील
~$\TMstroke$ चिन्हांच्या खंडांच्या क्रमिकेने
दर्शवलेले असेल. या क्रमिकेतील पहिला
!!{element} म्हणजे अवस्थांची संख्या आहे.
म्हणून पहिला $\TMstroke$ खंड सोडून इतर
सर्व पुसून टाकायचे आणि मग थांबायचे, एवढेच
उरते.

पुढील फलित विशेष उल्लेखनीय आहे:

\begin{thm}\ollabel{thm:universal-tm} असे एक \emph{सार्वत्रिक
  ट्यूरिंग यंत्र}~$U$ आहे की आदान $\tuple{e,n}$ वर
  सुरू केल्यास ते
  \begin{enumerate}
    \item $M_e$ हे आदान~$n$ वर थांबते तेव्हाच थांबते, आणि
    \item $M_e$ हे फलित $m$ देऊन थांबले, तर~$U$ देखील
    त्याच फलितासह थांबते.
  \end{enumerate}
  म्हणून $U$ हे $f\colon \Nat \times \Nat \pto \Nat$
  फलन संगणित करते; $M_e$ हे आदान~$n$ वर सुरू
  होऊन फलित~$m$ देत थांबले तर $f(e,n) = m$,
  आणि अन्यथा ते अपरिभाषित असते.
\end{thm}

\begin{proof}
  $U$ प्रत्यक्ष तयार करणे जवळजवळ अशक्य आहे, कारण
  ते अत्यंत गुंतागुंतीचे यंत्र आहे. पण ते कसे
  काम करते याची रूपरेषा देता येते आणि मग
  चर्च--ट्यूरिंग प्रबंधाचा आधार घेता येतो.
  सुरुवातीला $U$ च्या फितीवर $e$ $\TMstroke$
  चिन्हांचा खंड आणि त्यानंतर $n$~$\TMstroke$
  चिन्हांचा खंड असतो. ते प्रथम निर्देशांक~$e$
  आदान~$n$ च्या उजवीकडे “उलगडते.” यामुळे
  यंत्र~$M_e$ च्या सूचनांचे वर्णन करणाऱ्या
  संख्यांची यादी मिळते (म्हणजे~$\TMblank$
  चिन्हांनी वेगळे केलेले $\TMstroke$ चिन्हांचे
  खंड). त्यानंतर $U$ हे $M_e$ च्या वर्णनाच्या
  उजवीकडे $M_e$ च्या आरंभीच्या अवस्थेची
  संख्या आणि संख्या~$1$ लिहिते. (पुन्हा,
  या संख्या $\TMstroke$ चिन्हांच्या खंडांनी
  एकचिन्ही रूपात दर्शवलेल्या आहेत.) पुढे,
  आदानातील $n$~$\TMstroke$ चिन्हांचा खंड
  उजवीकडे नकलते---पण प्रत्येक $\TMstroke$
  चिन्हाच्या जागी तीन $\TMstroke$ चिन्हांचा
  खंड लिहिते ($\TMstroke$ या चिन्हाची संख्या~$3$
  आहे, तर~$\TMendtape$ ची संख्या $1$ आणि
  ~$\TMblank$ ची संख्या $2$ आहे, हे आठवा).
  अशा खंडांच्या क्रमिकेच्या डाव्या टोकाला
  ($U$ च्या फितीवर ते~$\TMblank$ चिन्हांनी
  वेगळे केलेले आहेत) ते एकच~$\TMstroke$
  चिन्ह लिहिते; हाच~$\TMendtape$ चा संकेतांक आहे.

  आता $U$ च्या फितीवर निर्देशांक~$e$, संख्या~$n$,
  आरंभीच्या अवस्थेचा संकेतांक (“सध्याची अवस्था”),
  आरंभीच्या शीर्षस्थानाची संख्या~$1$
  (“सध्याचे शीर्षस्थान”) आणि अनुकरणातील
  “फिती”वरील आरंभीचा मजकूर नोंदलेला असतो
  (म्हणजे~$M_e$ च्या “फिती”वरील “चिन्हांचे”
  संकेतांक दर्शवणाऱ्या~$\TMstroke$ चिन्हांच्या
  खंडांची, ~$\TMblank$ चिन्हांनी वेगळे केलेली
  क्रमिका).

  आदान~$n$ वर सुरू केलेले $M_e$ जे काम
  करेल त्याचे अनुकरण ते पुढीलप्रमाणे करते:
  \begin{enumerate}
    \item “सध्याच्या शीर्षस्थानाची” संख्या~$k$ शोधणे
    (सुरुवातीला ती~$1$ असते),
    \item “फितीवरील” $k$-व्या खंडाकडे जाऊन
    तेथील “चिन्ह” कोणते आहे ते पाहणे,
    \item\ollabel{find-inst}%
    सध्याची “अवस्था” आणि “चिन्ह” यांना जुळणारी
    सूचना शोधणे,
    \item “फितीवरील” $k$-व्या खंडाकडे
    परत जाऊन तेथील “चिन्हा”ऐवजी $M_e$ लिहील
    त्या चिन्हाचा संकेतांक लिहिणे,
    \item सध्याची “अवस्था” नोंदवलेली आहे
    तेथे शीर्ष नेऊन त्या संख्येऐवजी नव्या
    अवस्थेची संख्या लिहिणे,
    \item “फितीवरील स्थान” नोंदवलेले आहे
    तेथे जाऊन एक~$\TMstroke$ पुसणे किंवा
    एक~$\TMstroke$ जोडणे (सूचना अनुक्रमे
    डावीकडे किंवा उजवीकडे जायला सांगत असेल तर).
    \item पुन्हा हेच करणे.\footnote{येथे आपण काही सूक्ष्म
    अडचणींचा तपशील टाळत आहोत. उदाहरणार्थ,
    “सध्याचे शीर्षस्थान” मोजणाऱ्या संख्येत
    वाढ करताना $U$ ला अतिरिक्त जागा लागू शकते;
    अशा वेळी त्याला संपूर्ण “फीत” उजवीकडे
    हलवावी लागेल.}
  \end{enumerate}
  आदान~$n$ वर सुरू केलेले $M_e$ कधीच थांबत
  नसेल, तर $U$ देखील कधीच थांबत नाही;
  म्हणून त्याचे फलित अपरिभाषित असते.

  पायरी~\olref{find-inst} मध्ये $M_e$ च्या
  वर्णनात सध्याच्या “अवस्था”/“चिन्ह” जोडीला
  जुळणारी सूचना नसेल, तर $M_e$ थांबले असते.
  तसे झाल्यास $U$ आपल्या फितीवरील “फिती”च्या
  डावीकडचा भाग पुसते. $M_e$ च्या फितीवरील
  एक~$\TMstroke$ दर्शवणाऱ्या प्रत्येक तीन
  ~$\TMstroke$ चिन्हांच्या खंडासाठी ते
  स्वतःच्या फितीच्या डाव्या टोकाला एक
  $\TMstroke$ लिहिते आणि त्याच वेळी क्रमाने
  “फीत” पुसत जाते. हे संपल्यावर $U$ च्या
  फितीवर लांबी~$m$ असलेला $\TMstroke$
  चिन्हांचा एकच खंड उरतो.

  “फितीवर” तीन~$\TMstroke$ चिन्हांच्या
  खंडाऐवजी काही वेगळे आढळले, तर $U$
  लगेच थांबते. या बाबतीत $U$ च्या फितीवर
  $\TMstroke$ चिन्हांचा एकच खंड नसल्याने
  त्याचे फलित नैसर्गिक संख्या नसते; म्हणजे
  या बाबतीत $f(e,n)$ अपरिभाषित असते.
\end{proof}

\end{document}
